\documentclass[10pt]{article}

\usepackage[margin=1in]{geometry}
\usepackage{enumitem}
\usepackage{hyperref}
\usepackage{parskip}
\usepackage{titlesec}

\hypersetup{
    colorlinks=true,
    urlcolor=blue,
    linkcolor=black
}

\setlist[itemize]{leftmargin=1.5em}

\titleformat{\section}
{\large\bfseries}
{}{0pt}{}

\begin{document}
\begin{center}
    {\LARGE\bfseries GEOMETRY (REGENTS) SUMMER COURSE SYLLABUS}
\end{center}

\section*{Course Information}

\begin{tabular}{ll}
\textbf{Course:} & Geometry (Regents) Summer School \\
\textbf{Teacher:} & Mr. Ben Huang \\
\textbf{Email:} & \href{mailto:bhuang@windsorschool.com}{bhuang@windsorschool.com}
\end{tabular}

\section*{Welcome}

Welcome to Summer Geometry! My name is Mr. Ben Huang, and I will be your Geometry teacher this summer.

I was born in China and speak both Cantonese and Mandarin. I have been teaching mathematics at both the college and high school levels, either full-time or part-time, since 2010. In addition to classroom teaching, I tutor students preparing for the SAT, AP Mathematics, and New York State Regents examinations.

In 2025, I moved from Central Pennsylvania to New York to join my wife and our then newly born daughter. I am excited to be part of the New York educational community and look forward to working with students here.

Outside of teaching, I enjoy solving challenging mathematics and programming problems, as well as developing web applications and video games. One of my ongoing game development projects can be found here:

\medskip

\noindent
\url{https://benhuangbmj.itch.io/genssome-akademy}

\medskip

I look forward to working with you this summer and helping you develop both your mathematical understanding and your confidence as mathematical thinkers.

\section*{Course Description}

This accelerated six-week course is aligned with the New York State Next Generation Mathematics Learning Standards and prepares students for the New York State Geometry Regents Examination. Students will develop geometric reasoning, construct mathematical proofs, analyze figures using transformations and coordinate methods, and apply geometry to solve real-world problems.

\section*{Learning Objectives}

By the end of this course, students will be able to:

\begin{itemize}
    \item Analyze and describe geometric relationships.
    \item Construct and justify mathematical proofs.
    \item Apply transformations to establish congruence and similarity.
    \item Solve problems involving triangles, polygons, and circles.
    \item Use coordinate geometry to investigate geometric properties.
    \item Apply right triangle trigonometry.
    \item Calculate perimeter, area, surface area, and volume.
    \item Apply geometry to solve real-world problems.
    \item Communicate mathematical reasoning clearly and precisely.
\end{itemize}

\section*{Major Course Topics}

\begin{itemize}
    \item Foundations of Geometry
    \item Geometric Constructions
    \item Transformations
    \item Congruence and Proof
    \item Parallel Lines and Polygons
    \item Similarity
    \item Right Triangle Trigonometry
    \item Circles
    \item Coordinate Geometry
    \item Geometric Measurement
    \item Surface Area and Volume
    \item Regents Examination Review
\end{itemize}

\section*{Required Textbook}

\textbf{Geometry (NY Next Generation Edition)}

Publisher: Perfection Learning

ISBN: 9781690305200

Additional notes, assignments, and review materials will be posted on Google Classroom throughout the course.

\section*{Materials}

Students should bring the following materials to class every day:

\begin{itemize}
    \item Geometry (NY Next Generation Edition) textbook
    \item TI-83 or TI-84 graphing calculator
    \item Notebook or binder
    \item Pencil and eraser
    \item Ruler (straightedge)
    \item Compass
    \item Protractor
\end{itemize}

\section*{How We Will Learn}

This class is designed around the ideas of thinking, collaboration, and problem solving. Rather than simply watching examples and copying procedures, you will regularly work on meaningful mathematical tasks that encourage you to reason, discuss ideas, and learn from one another.

To support this learning environment, you can expect the following classroom practices:

\begin{itemize}
    \item \textbf{Randomized groups and seating.} You will frequently work with different classmates. This allows everyone to learn from a variety of perspectives and develop the ability to collaborate effectively with others.

    \item \textbf{Non-permanent vertical workspaces.} Many activities will be completed while standing and working on erasable vertical surfaces. These workspaces encourage experimentation, discussion, revision, and the sharing of mathematical ideas.

    \item \textbf{Thinking before telling.} You will often be asked to make sense of a problem before receiving direct instruction. Productive struggle is an important part of learning mathematics, and making mistakes are a natural part of the learning process.

    \item \textbf{Discussion and justification.} You are expected to explain your reasoning, ask questions, compare different solution strategies, and communicate your mathematical thinking clearly.

    \item \textbf{Active participation.} Mathematics is learned by doing. Every student is expected to engage in problem solving, contribute to group discussions, and support the learning of classmates.
\end{itemize}

Our classroom is a place where curiosity, perseverance, respectful collaboration, and thoughtful risk-taking are valued. The goal is not simply to arrive at the correct answer, but to become a stronger mathematical thinker.

\section*{Attendance}

Because this is an accelerated six-week summer course, attendance is essential.

Students are expected to:

\begin{itemize}
    \item Arrive to class on time every day.
    \item Remain in class for the entire three-hour session.
    \item Attend every class whenever possible.
    \item Complete any missed work promptly if an absence is unavoidable.
\end{itemize}

Missing even one class can make it difficult to keep up with the pace of the course.

\section*{Assessment and Grading}

Your course grade will be determined as follows:

\begin{center}
\begin{tabular}{lr}
Assessments (quizzes, exams, and major projects) & 50\% \\
Homework & 25\% \\
Midterm/Final Exam & 25\%
\end{tabular}
\end{center}

Homework will be assigned every day and should be completed on time. Multiple quizzes and assessments will be administered throughout the course to monitor your progress and prepare you for the Geometry Regents Examination.

\section*{Classroom Expectations}

Students are expected to:

\begin{itemize}
    \item Come prepared and ready to learn each day.
    \item Participate actively in class activities and discussions.
    \item Complete all assignments on time.
    \item Show mathematical reasoning and justify solutions when appropriate.
    \item Respect classmates, teachers, and school property.
    \item Demonstrate academic honesty on all assignments and assessments.
\end{itemize}

\section*{Google Classroom}

Google Classroom will be our primary online learning platform. Students should join the class as soon as possible and check it daily for:

\begin{itemize}
    \item Homework assignments
    \item Class notes
    \item Announcements
    \item Review materials
    \item Important course updates
\end{itemize}

\section*{Regents Examination}

Students enrolled in this course are expected to take the New York State Geometry Regents Examination at the conclusion of the course.

\end{document}